% ------------------------------------------------------------------ title page
\begin{titlepage}
\thispagestyle{empty}
\begin{tikzpicture}[remember picture, overlay]
  % Orange band across the top of the page (text area starts 2.6cm below the page edge).
  \fill[ttorange] (current page.north west) rectangle
    ($(current page.north east)+(0,-10.4cm)$);
  \fill[ttdeep] ($(current page.north west)+(0,-10.4cm)$) rectangle
    ($(current page.north east)+(0,-10.68cm)$);

  % A decorative automata grid: the learned state of a small machine, rendered in the band.
  \begin{scope}[shift={($(current page.north east)+(-4.75cm,-1.35cm)$)}, scale=0.52,
                every node/.style={inner sep=0pt}]
    \foreach \r in {0,...,7}{
      \foreach \c in {0,...,13}{
        \pgfmathsetmacro{\v}{mod(\r*7+\c*5+int(\r*\c/2),9)}
        \ifdim \v pt > 5.5pt
          \fill[white, rounded corners=0.4pt] (\c,-\r) rectangle ++(0.8,0.8);
        \else
          \ifdim \v pt > 3.0pt
            \fill[white, opacity=0.40, rounded corners=0.4pt] (\c,-\r) rectangle ++(0.8,0.8);
          \else
            \draw[white, opacity=0.28, line width=0.5pt, rounded corners=0.4pt]
              (\c,-\r) rectangle ++(0.8,0.8);
          \fi
        \fi
      }
    }
  \end{scope}
\end{tikzpicture}

\vspace*{0.2cm}
{\sffamily\color{white}
  {\fontsize{12}{14}\selectfont\bfseries WHITE PAPER}\par
  \vspace{0.18cm}
  {\fontsize{10}{12}\selectfont\color{ttsand} Version \ttversion\ \ \textbar\ \ \ttdate}\par
  \vspace{1.35cm}
  {\fontsize{40}{44}\selectfont\bfseries torchtsetlin\par}
  \vspace{0.42cm}
  {\fontsize{15.5}{20}\selectfont
   A GPU-native, PyTorch-idiomatic framework\\ for Tsetlin machines\par}
}

\vspace{1.5cm}

{\sffamily
 {\large\bfseries\color{ttink} Vajira Thambawita}\par
 \vspace{0.12em}
 {\color{ttgray} Simula Metropolitan Center for Digital Engineering, Oslo, Norway}\par
 {\color{ttgray} \texttt{vajira@simula.no}}\par
}

\vspace{0.85cm}

\begin{tikzpicture}
  \node[inner sep=0pt] (abs) {%
    \begin{minipage}{0.925\linewidth}
      \small
      {\sffamily\bfseries\color{ttorange}Abstract.}\enspace
      The Tsetlin machine is a propositional-logic learner: it builds classifiers out of
      conjunctive clauses over Boolean literals, updated not by gradient descent but by
      reward\slash penalty feedback to teams of Tsetlin automata. Its models are legible by
      construction and its inference is bit-level cheap, which is why it has become
      attractive for clinical decision support, edge inference and hardware accelerators.
      Its software, however, has grown up outside the deep-learning ecosystem: the reference
      implementations are C and CUDA kernels behind scikit-learn-shaped wrappers, each model
      family a separate package, with no shared tensor abstraction, no common training loop
      and no route into PyTorch tooling.
      \par\smallskip
      \ttlibname closes that gap. It expresses the whole Tsetlin machine family --- vanilla,
      weighted, coalesced, regression and convolutional --- as ordinary
      \code{torch.nn.Module} objects whose learnable state is a tensor of automaton counters.
      A model therefore moves across devices with \code{.to(device)}, serialises with
      \code{state\_dict()}, and learns from a mini-batch with a single \code{model.update(x,
      y)} call that plays the role of \code{loss.backward(); optimizer.step()}. The central
      technical contribution is a \emph{batched} formulation of Type~I\slash Type~II feedback
      in which per-example stochastic automaton transitions are replaced by event counts
      accumulated with two matrix products and realised as binomial draws --- one commit per
      batch instead of one per example. On an RTX~3090 this reaches
      $2.2\times10^{5}$~training examples\slash s and up to $273\times$ the throughput of the
      same code on a 16-core CPU, while matching the sequential algorithm's accuracy at
      moderate batch sizes. Around the models the package ships Booleanization encoders, a
      Keras-style trainer with callbacks, metrics, closed-form global and local
      interpretability, and visualisation --- everything a Tsetlin machine workflow needs, in
      one PyTorch-shaped package.
    \end{minipage}};
  \begin{scope}[on background layer]
    \draw[ttorange!35, line width=0.8pt, rounded corners=4pt, fill=ttsand]
      ($(abs.north west)+(-1.1em,0.95em)$) rectangle ($(abs.south east)+(1.1em,-0.95em)$);
  \end{scope}
\end{tikzpicture}

\vfill

{\sffamily\scriptsize\color{ttgray}
 \setlength{\parskip}{0.15em}
 \textcolor{ttorange}{\rule[0.25em]{2.2em}{1.6pt}}\par
 \textbf{Software:} \ttlibname \ttversion\ \ \textbar\ \ MIT licence \ \textbar\ \
 \texttt{pip install torchtsetlin}\par
 \textbf{Code:} \url{https://github.com/vlbthambawita/torchtsetlin} \ \ \textbar\ \
 \textbf{Docs:} \url{https://vlbthambawita.github.io/torchtsetlin/}\par
 \textbf{Measurements} in Section~\ref{sec:benchmarks} were produced by the benchmark driver
 shipped with the repository (\texttt{benchmarks/run\_benchmarks.sh}) and can be reproduced on
 other hardware with one command.\par
}
\end{titlepage}
